% Generated draft: requires mathematical and source review.
\documentclass[11pt]{article}
\usepackage[margin=25mm]{geometry}
\usepackage{fontspec}
\setmainfont{Latin Modern Roman}
\usepackage{amsmath,amssymb,mathtools,unicode-math}
\setmathfont{Latin Modern Math}
\usepackage{enumitem,needspace,xurl,hyperref}
\setlength{\parindent}{0pt}
\setlength{\parskip}{5pt}
\setlength{\emergencystretch}{3em}
\newcommand{\N}{\mathbb{N}}
\newcommand{\Z}{\mathbb{Z}}
\newcommand{\Q}{\mathbb{Q}}
\newcommand{\R}{\mathbb{R}}
\newcommand{\C}{\mathbb{C}}
\newcommand{\F}{\mathbb{F}}
\newcommand{\A}{\mathbb{A}}
\newcommand{\PP}{\mathbb{P}}
\newcommand{\E}{\mathbb{E}}
\newcommand{\eps}{\varepsilon}
\newcommand{\dd}{\,\mathrm{d}}
\begin{document}
\section*{UnsolvedMath ID: 800138}
\textbf{Source problem code:} catalogue entry
\subsection{Definitions and mathematical statement}
Let \(\mathbb K\) be a number field, that is, a finite extension of \(\Q\). The supplied problem asks about a quantity denoted by \(\Delta_{\mathbb K}(x)\), but does not define this quantity, its normalization, or the domain of \(x\). Consequently, the following is a conditional formulation of the requested task, pending those definitions.

Write \(D_{\mathbb K}\subseteq\R\) for the intended, unspecified domain, and suppose that the intended function is \(\Delta_{\mathbb K}:D_{\mathbb K}\to\R\). Define its negative set by
\[
N_{\mathbb K}=\{x\in D_{\mathbb K}:\Delta_{\mathbb K}(x)<0\}.
\]
For each of
\[
\mathbb K=\Q(\sqrt{5})\qquad\text{and}\qquad\mathbb K=\Q(\sqrt{13}),
\]
localize the first negative occurrence of \(\Delta_{\mathbb K}(x)\); the source also permits investigating another simple number field, without specifying what constitutes a simple field. If ``first'' means a least argument, the target is \(\min N_{\mathbb K}\), provided that this minimum exists. The source does not specify whether an infimum or an interval containing the first negative occurrence is intended when there is no least argument.

More generally, obtain an upper bound for the first sign change of \(\Delta_{\mathbb K}(x)\) as \(\mathbb K\) ranges over a family of number fields. The family, the parameters controlling the bound, the allowed hypotheses, the initial sign and starting range, and the convention for a sign change are unspecified. A first negative occurrence and a first sign change therefore cannot be identified without additional information.

The accompanying comment, attributed to Alexandre Bailleul (Universit\'e Paris--Saclay), concerns highly biased, appropriately normalized races between prime quadratic residues and prime quadratic nonresidues for composite moduli \(q\). For a positive composite integer \(q\), let
\[
G_q=(\Z/q\Z)^\times,\qquad
G_q^2=\{g^2:g\in G_q\},\qquad
\rho(q)=[G_q:G_q^2],\qquad
\operatorname{rad}(q)=\prod_{p\mid q}p,
\]
where the product is over the distinct prime divisors of \(q\). The comment attributes to Fiorilli [Fio] a conjecture that the Skewes number for this race, described as the smallest real argument at which the counting function for prime quadratic residues overcomes the bias, is roughly of the form
\[
\exp\!\left(\exp\!\left(\frac{\rho(q)}{\log\operatorname{rad}(q)}\right)\right).
\]
The normalization, precise crossing inequality, hypotheses, range of moduli, and mathematical meaning of ``roughly'' are not supplied. This contextual formula is not a specified bound for the number-field problem. The comment also states that Fiorilli and Jouve [FJ] and Bailleul [Bai] investigated similar highly biased races in number fields.
\subsection{Short English statement}
Locate the first negative occurrence of the source's quantity \(\Delta_{\mathbb K}(x)\) for \(\Q(\sqrt{5})\), \(\Q(\sqrt{13})\), or another simple number field, and seek upper bounds for its first sign change in a family of number fields. The quantity itself and several conventions are missing, so this remains a conditional statement requiring clarification. The supplied registry labels the problem open; this label is not an independently verified assessment.
\subsection{Sources}
\begin{itemize}
\item Local UnsolvedMath registry, supplied complete source record, permanent ID \(800138\), exact source problem code \texttt{catalogue entry}, title ``Comparative Prime Number Theory --- Problem 7.'' The record attributes the problem to Hamieh, Kadiri, Martin, and Ng, with proposed year 2024, and supplies the status \texttt{open}.
\item Hamieh, Kadiri, Martin, and Ng, \emph{Comparative Prime Number Theory Problem List} (2024), Problem (\#7), \url{https://arxiv.org/abs/2407.03530}. The record reports access on 2026-10-16 and extraction by \texttt{arxiv\_latex}; these are supplied provenance claims, not independently verified access or source verification.
\item The embedded comment is attributed to Alexandre Bailleul, Universit\'e Paris--Saclay. It attributes the contextual Skewes-number conjecture to Fiorilli [Fio] and mentions related number-field investigations by Fiorilli and Jouve [FJ] and Bailleul [Bai]. The supplied record gives no full bibliographic entries or links for these three reference labels.
\end{itemize}
\end{document}
